% !TEX root = ../report.tex

\chapter{Исследовательская часть}

В данном разделе исследуется зависимость производительности разработанной системы от двух факторов: наличия и типа индексов в PostgreSQL и наличия внешнего кэша Redis. Цель --- количественно оценить, как эти решения влияют на время выполнения запросов и время отклика приложения, и определить условия, при которых они дают эффект.

Проводится исследование, а не эксперимент: измеряются свойства уже реализованной системы на воспроизводимом наборе данных, без построения формальной модели и проверки статистических гипотез.

\section{Методика измерений}

Измерения выполняются на изолированной базе данных, наполненной синтетическим набором данных. Набор формируется набором миграций и скриптом генерации; объём приведён в таблице~\ref{tbl:research_dataset}. Имена товаров составляются из словаря частотных токенов и редкого токена, что задаёт управляемую селективность для запросов текстового поиска.

\begin{table}[H]
    \caption{Параметры набора данных основного прогона}
    \label{tbl:research_dataset}
    \begin{tabular}{|p{0.5\textwidth}|p{0.42\textwidth}|}
        \hline
        \textbf{Параметр} & \textbf{Значение} \\ \hline
        Товаров (синтетических) & 40\,000 \\ \hline
        Заказов & 80\,000 \\ \hline
        Отзывов & 50\,000 \\ \hline
        Концентрация заказов & первые 100 покупателей \\ \hline
    \end{tabular}
\end{table}

Используются два инструмента измерения:

\begin{itemize}
    \item для запросов к базе данных --- команда \texttt{EXPLAIN (ANALYZE, BUFFERS, FORMAT JSON)}, из плана извлекаются время выполнения, время планирования, тип корневого узла плана и наличие узла сортировки;
    \item для интерфейса доступа --- нагрузочный тест: фиксированное число одновременных HTTP-запросов к маршрутам \texttt{/api/v1/...}, с измерением среднего времени ответа и перцентилей \texttt{P95}, \texttt{P99}.
\end{itemize}

Каждый запрос к базе данных измеряется в нескольких повторах; первый прогон отбрасывается как разогрев плана и буферного кэша, по остальным вычисляются среднее и среднеквадратическое отклонение (СКО). Перед каждой серией выполняется \texttt{ANALYZE} для актуализации статистики планировщика.

\section{Исследование 1: наличие и тип индексов}

\subsection{Наличие вторичных индексов}

В первом подыспытании сравниваются два варианта схемы: без вторичных индексов и с индексами, создаваемыми миграциями. Набор запросов покрывает выборку каталога по диапазону цены, выборку заказов покупателя, выборку отзывов товара и хранимую функцию статистики. Результаты приведены в таблице~\ref{tbl:research_presence} и на рисунке~\ref{fig:research_presence}.

\begin{figure}[H]
    \centering
    \includegraphics[width=0.9\textwidth]{research_presence.png}
    \caption{Время выполнения запросов при наличии и отсутствии индексов}
    \label{fig:research_presence}
\end{figure}

Для запросов с выборкой по индексируемому столбцу наличие индекса сокращает время выполнения на один--два порядка: планировщик переходит от последовательного сканирования (\texttt{Seq Scan}) к сканированию по индексу. Для хранимой функции статистики индексы из набора миграций не дают выигрыша, а в отдельных прогонах время незначительно возрастает: запрос соединяет три таблицы с агрегацией, и доступные индексы не покрывают его план. Это показывает, что индекс полезен не всегда и должен соответствовать конкретному запросу.

\subsection{Тип индекса: простой и составной}

Во втором подыспытании сравниваются простой индекс \texttt{(user\_id)} и составной \texttt{(user\_id, created\_at desc)} для запроса выборки заказов покупателя с сортировкой по дате. Результаты --- в таблице~\ref{tbl:research_type}.

При простом индексе план содержит узел сортировки (\texttt{Sort}): индекс отбирает строки покупателя, после чего они сортируются по дате. Составной индекс хранит строки уже в нужном порядке, поэтому узел сортировки из плана исчезает, а время выполнения снижается. Эффект проявляется при достаточном числе заказов у покупателя; на равномерном распределении он незаметен, поэтому заказы в наборе сконцентрированы на части покупателей.

\subsection{Тип индекса: B-tree и GIN для текстового поиска}

Третье подыспытание сравнивает два типа индекса для подстрочного поиска по имени товара (\texttt{name ILIKE '\%токен\%'}). Индекс B-tree для шаблона с ведущим символом подстановки неприменим, поэтому выполняется последовательное сканирование. Индекс GIN с классом операторов \texttt{gin\_trgm\_ops} (расширение \texttt{pg\_trgm}) индексирует триграммы и обслуживает такой поиск через \texttt{Bitmap Heap Scan}. Результаты --- в таблице~\ref{tbl:research_type}.

При низкой селективности запроса GIN-индекс сокращает время выполнения более чем на порядок относительно последовательного сканирования. Рисунок~\ref{fig:research_scaling} показывает, что разрыв растёт с увеличением числа товаров: время последовательного сканирования растёт линейно по размеру таблицы, а время поиска по GIN-индексу остаётся практически постоянным.

\begin{figure}[H]
    \centering
    \includegraphics[width=0.9\textwidth]{research_scaling.png}
    \caption{Время текстового поиска в зависимости от размера набора: Seq Scan и GIN}
    \label{fig:research_scaling}
\end{figure}

\section{Исследование 2: влияние Redis-кэширования}

Во втором исследовании оценивается влияние внешнего кэша Redis на время отклика приложения. Нагрузка подаётся на четыре маршрута чтения: карточка товара, каталог, отзывы товара и список категорий. Сравниваются два сценария: Redis отключён и Redis включён с прогретым кэшем. Результаты приведены в таблице~\ref{tbl:research_http} и на рисунке~\ref{fig:research_http}; доля попаданий в кэш и относительное число обращений к PostgreSQL --- в таблице~\ref{tbl:research_cache_meta}.

\begin{figure}[H]
    \centering
    \includegraphics[width=0.9\textwidth]{research_http.png}
    \caption{Среднее время ответа по эндпоинтам: Redis отключён и включён}
    \label{fig:research_http}
\end{figure}

Включение кэша снижает среднее время ответа и, в большей степени, хвостовые перцентили \texttt{P95} и \texttt{P99}: повторные чтения обслуживаются из памяти, без обращения к PostgreSQL. При прогретом кэше доля попаданий близка к единице, а число строк, читаемых из PostgreSQL, снижается практически до нуля. Величина выигрыша зависит от маршрута (см. таблицу~\ref{tbl:research_http}): для часто запрашиваемой карточки товара относительное снижение времени наибольшее, тогда как для и без того дешёвого списка категорий эффект мал.

% Generated by scripts/research/render_results.py
\begin{table}[H]
\centering
\caption{Среднее время выполнения запросов при различном наборе индексов}
\label{tbl:research_presence}
\begin{tabular}{|p{0.32\textwidth}|p{0.40\textwidth}|r|}
\hline
Запрос & Вариант & Время, мс \\ \hline
catalog\_price\_range & без вторичных индексов & 9.351 \\ \hline
catalog\_price\_range & с B-tree индексами миграций & 0.073 \\ \hline
orders\_by\_user & без вторичных индексов & 7.448 \\ \hline
orders\_by\_user & с B-tree индексами миграций & 1.977 \\ \hline
pending\_orders\_by\_date & без вторичных индексов & 11.252 \\ \hline
pending\_orders\_by\_date & с B-tree индексами миграций & 0.394 \\ \hline
reviews\_by\_product & без вторичных индексов & 1.156 \\ \hline
reviews\_by\_product & с B-tree индексами миграций & 0.288 \\ \hline
seller\_statistics\_function & без вторичных индексов & 11.983 \\ \hline
seller\_statistics\_function & с B-tree индексами миграций & 12.894 \\ \hline
\end{tabular}
\end{table}

\begin{table}[H]
\centering
\caption{Влияние типа индекса на время выполнения}
\label{tbl:research_type}
\begin{tabular}{|p{0.42\textwidth}|c|p{0.22\textwidth}|r|}
\hline
Случай & Узел Sort & Корневой узел & Время, мс \\ \hline
orders\_by\_user: простой индекс (user\_id) & yes & Limit & 2.100 \\ \hline
orders\_by\_user: составной индекс (user\_id, created\_at desc) & no & Limit & 0.528 \\ \hline
text\_search: только B-tree (для '\%x\%' неприменим -> Seq Scan) & no & Seq Scan & 21.962 \\ \hline
text\_search: GIN + pg\_trgm & no & Bitmap Heap Scan & 0.442 \\ \hline
\end{tabular}
\end{table}

\begin{table}[H]
\centering
\caption{Время текстового поиска в зависимости от размера набора}
\label{tbl:research_scaling}
\begin{tabular}{|r|p{0.26\textwidth}|r|p{0.22\textwidth}|}
\hline
Размер & Вариант & Время, мс & Узел \\ \hline
10000 & Seq Scan (без GIN) & 9.234 & Seq Scan \\ \hline
10000 & GIN + pg\_trgm & 0.343 & Bitmap Heap Scan \\ \hline
100000 & Seq Scan (без GIN) & 47.672 & Seq Scan \\ \hline
100000 & GIN + pg\_trgm & 0.843 & Bitmap Heap Scan \\ \hline
500000 & Seq Scan (без GIN) & 94.959 & Gather \\ \hline
500000 & GIN + pg\_trgm & 2.148 & Bitmap Heap Scan \\ \hline
\end{tabular}
\end{table}

\begin{table}[H]
\centering
\caption{Нагрузочное тестирование: время ответа по эндпоинтам (влияние Redis)}
\label{tbl:research_http}
\begin{tabular}{|p{0.20\textwidth}|p{0.28\textwidth}|r|r|r|}
\hline
Эндпоинт & Сценарий & RPS & Среднее, мс & P99, мс \\ \hline
Карточка товара & Redis отключён & 4999.96 & 2.256 & 18.156 \\ \hline
Каталог & Redis отключён & 3893.96 & 2.935 & 11.712 \\ \hline
Отзывы & Redis отключён & 3977.65 & 2.879 & 5.590 \\ \hline
Категории & Redis отключён & 4596.96 & 2.417 & 4.439 \\ \hline
Карточка товара & Redis включён, прогретый кэш & 11019.35 & 0.965 & 2.638 \\ \hline
Каталог & Redis включён, прогретый кэш & 7135.62 & 1.561 & 3.389 \\ \hline
Отзывы & Redis включён, прогретый кэш & 9186.10 & 1.198 & 2.559 \\ \hline
Категории & Redis включён, прогретый кэш & 6872.77 & 1.570 & 3.200 \\ \hline
\end{tabular}
\end{table}

\begin{table}[H]
\centering
\caption{Доля попаданий в кэш и обращения к PostgreSQL}
\label{tbl:research_cache_meta}
\begin{tabular}{|p{0.34\textwidth}|r|r|r|}
\hline
Сценарий & Попаданий & Промахов & Доля попаданий \\ \hline
Redis отключён & 0 & 0 & 0.000 \\ \hline
Redis включён, прогретый кэш & 1500 & 0 & 1.000 \\ \hline
\end{tabular}
\end{table}


\section*{Вывод}
\addcontentsline{toc}{section}{Вывод}

Проведённое исследование показало, что индекс ускоряет выборку только при соответствии его структуры запросу: наличие индекса по фильтруемому столбцу сокращает время на один--два порядка, составной индекс дополнительно устраняет сортировку, а для подстрочного поиска требуется индекс типа GIN, тогда как B-tree неприменим. Внешний кэш Redis снижает время отклика приложения и особенно хвостовые задержки за счёт сокращения обращений к PostgreSQL, причём сильнее всего для дорогих запросов. PostgreSQL остаётся источником истины, а индексы и кэш выступают средствами оптимизации производительности в рамках одной и той же схемы данных.
